\section{Initial Algebras}

\begin{defns}\ 
    \begin{enumerate}[-]
            \item An \textsl{$\fun L_X$-algebra} is a type $C$ provided with an element $c_{\fun{nil}}\colon C$ and a function $c_{\fun{cons}}\colon X\to C\to C$. The type of such algebras is
            \[
                \fun L_X\fun{Alg}\defeq\sum_{C\colon\univ}C\times(X\to C\to C).
            \]
            \item Let $\mathcal C\defeq\dep{C,(c_{\fun{nil}},c_{\fun{cons}})}$ and $\mathcal D\defeq\dep{D,(d_{\fun{nil}}, d_{\fun{cons}})}$ be two $\fun L_X$-algebras. An \textsl{$\fun L_X$-morphism} between them is a map $h\colon C\to D$ such that $h(c_{\fun{nil}})\eq Dd_{\fun{nil}}$ and $h(c_{\fun{cons}}(x,c))\eq Dd_{\fun{cons}}(x,h(c))$ for all $x\colon X$ and $c\colon C$. The type of such morphisms is
        \[
            \fun L_X\Hom(\mathcal C,\mathcal D)\defeq
            \sum_{h}
                \bigl(
                    h(c_{\fun{nil}})
                    \eq D
                    d_{\fun{nil}}
                \bigr)
                \times
                \prod_{x,c}
                    h(c_{\fun{cons}}(x,c))
                    \eq D
                    d_{\fun{cons}}(x, h(c)).
        \]
    
        \item An $\fun L_X$-algebra $\mathcal I$ is \textsl{homotopy-initial} (a.k.a.~\textsl{h-initial}) if for any $\fun L_X$-algebra $\mathcal C$ the type of $\fun L_X$-morphisms from $\mathcal I$ to $\mathcal C$ is contractible, i.e.,
        \[
            \fun{isHinit}_{\fun L_X}(\mathcal I)
                \defeq
                \prod_{\mathcal C\colon\fun L_X\fun{Alg}}
                    \fun{isContr}(
                        \fun L_X\Hom(\mathcal I,\mathcal C)
                    )
        \]
        is inhabited.
    \end{enumerate}
\end{defns}

\begin{lem}\label{lem:L_X-morphisms-are-morphisms}
    Given $\fun L_X$-algebras $\mathcal{A}\defeq\dep{A,(a_{\fun{nil}}, a_{\fun{cons}})}$, $\mathcal{B}\defeq\dep{B,(b_{\fun{nil}},b_{\fun{cons}})}$, and $\mathcal{C}\defeq\dep{C,(c_{\fun{nil}}, c_{\fun{cons}})}$, the following hold:
    \begin{enumerate}[a),font=\upshape]
        \item The identity function $\id_A\colon A\to A$ is an $\fun L_X$-morphism from $\mathcal{A}$ to $\mathcal{A}$.
        \item If $f\colon A\to B$ and $g\colon B\to C$ are $\fun L_X$-morphisms, then the composition $g\circ f\colon A\to C$ is an $\fun L_X$-morphism from $\mathcal{A}$ to $\mathcal{C}$.
        \item If $f\colon A\to B$ is an $\fun L_X$-morphism and $g\colon B\to A$ is a quasi-inverse of $f$, then the quasi-inverse $g$ is an $\fun L_X$-morphism from $\mathcal{B}$ to $\mathcal{A}$.
    \end{enumerate}
\end{lem}

\begin{proof}\ 
    \begin{enumerate}[a)]
        \item We need to show that $\id_A(a_{\fun{nil}})
        \eq{}a_{\fun{nil}}$ and $\id_A(a_{\fun{cons}}(x,a))\eq{}a_{\fun{cons}}(x,\id_A(a))$ for all $x\colon X$ and $a\colon A$. Both equalities hold definitionally.
        
        \item For the base case, we have
        \[
            (g\circ f)(a_{\fun{nil}})
                \equiv g(f(a_{\fun{nil}}))
                \eq Cg(b_{\fun{nil}})
                \eq Cc_{\fun{nil}}.
        \]
        For the inductive step, let $x\colon X$ and $a\colon A$. Then,
        \begin{align*}
            (g\circ f)(a_{\fun{cons}}(x,a))
                &\equiv g(f(a_{\fun{cons}}(x,a)))\\
                &\eq Cg(b_{\fun{cons}}(x,f(a)))\\
                &\eq Cc_{\fun{cons}}(x,g(f(a)))\\
                &\equiv c_{\fun{cons}}(x,(g\circ f)(a)).
        \end{align*}
        
        \item We have homotopies $\eta\colon f\circ g\sim\id_B$ and $\vartheta\colon g\circ f\sim\id_A$. We must show that $g$ preserves the algebra constructors. For the base case, we apply $\fun{ap}_g$ to the equality $f(a_{\fun{nil}})\eq Bb_{\fun{nil}}$ to obtain $g(f(a_{\fun{nil}}))\eq Ag(b_{\fun{nil}})$. By the homotopy $\vartheta$, we know $g(f(a_{\fun{nil}}))\eq Aa_{\fun{nil}}$. Transitivity of equality thus yields $g(b_{\fun{nil}})\eq Aa_{\fun{nil}}$.
        
        For the inductive step, let $x\colon X$ and $b\colon B$. Since $f$ is an $\fun L_X$-morphism, we have $f(a_{\fun{cons}}(x,g(b)))\eq Bb_{\fun{cons}}(x, f(g(b)))$. The homotopy $\eta$ implies $f(g(b))\eq Bb$. Then substitution gives
        \[
            f(a_{\fun{cons}}(x,g(b)))
            \eq Bb_{\fun{cons}}(x,b).
        \]
        Applying $g$ to both sides yields
        \[
            g(f(a_{\fun{cons}}(x,g(b))))
            \eq Ag(b_{\fun{cons}}(x,b)).
        \]
        By the homotopy $\vartheta$, the \lhs is equal to $a_{\fun{cons}}(x,g(b))$. Therefore, we conclude
        \[
            a_{\fun{cons}}(x,g(b))
            \eq Ag(b_{\fun{cons}}(x,b)),
        \]
        which shows that $g$ is an $\fun L_X$-morphism.\qedhere
    \end{enumerate}
\end{proof}

\begin{thm}
    The type of $h$-initial $\fun L_X$-algebras is a mere proposition.
\end{thm}

\begin{proof}
    Let $\mathcal I$ and $\mathcal J$ be two $h$-initial $\fun L_X$-algebras. Then, $\fun L_X\Hom(\mathcal I,\mathcal J)$ is contractible, hence inhabited by, say, $f\colon\mathcal I\to\mathcal J$. By symmetry, there is also $g\colon\mathcal J\to\mathcal I$. Since  Lemma~\ref{lem:L_X-morphisms-are-morphisms} implies that $g\circ f$ and $\id_{\mathcal I}$ inhabit $\fun L_X\Hom(\mathcal I,\mathcal I)$, which is contractible, $g\circ f\eq{\fun L_X\Hom(\mathcal I,\mathcal I)}\id_{\mathcal I}$. The same reasoning proves that $f\circ g\eq{\fun L_X\Hom(\mathcal J,\mathcal J)}\id_{\mathcal J}$.

    It follows that $f$ and $g$ are quasi-inverses and so $\mathcal I\simeq\mathcal J$. The conclusion follows by the univalence axiom.
\end{proof}

\begin{thm}
    The $\fun L_X$-algebra $\dep{\fun L_X,(\fun{nil},\fun{cons})}$ is homotopy initial.
\end{thm}

\begin{proof}
    Let $\mathcal C\defeq\dep{C,(c_{\fun{nil}},c_{\fun{cons}})}$ be an $\fun L_X$-algebra.

    Define $u\colon\fun L_X\to C$ inductively by
    \begin{itemize}
        \item $u(\fun{nil})\defeq c_{\fun{nil}}$,
        \item $u(\fun{cons}(a,\ell))\defeq c_{\fun{cons}}(a,u(\ell))$.
    \end{itemize}
    By definition, $u$ is an $\fun L_X$-morphism.

    Given any other $\fun L_X$-morphism $v\colon\fun L_X\to C$, we have
    \begin{itemize}
        \item $v(\fun{nil})\eq{C} c_{\fun{nil}}$,
        \item $v(\fun{cons}(a,\ell))\eq{C} c_{\fun{cons}}(a,v(\ell))$.
    \end{itemize}
    Therefore, by Remark~\ref{rem:List-induction-uniqueness}, we deduce that $u\sim v$, hence $u\eq{\fun L_X\Hom(\fun L_X,C)}v$. Consequently, $\fun L_X\Hom(\fun L_X,C)$ is contractible with center $u$.
    %
    
\end{proof}

\begin{cor}
    The $\fun L_{\fun1}$-algebra $\dep{\N,(0, \fun{succ})}$ is homotopy initial.
\end{cor}

\begin{proof}
    By the theorem, the $\fun L_{\fun1}$-algebra $\dep{\fun L_{\fun1},(\fun{nil}, \fun{cons})}$ is $h$-initial. Since we have established in Example~\ref{xmpl:N-from-List(1)} the equality $\N\eq{\univ}\fun L_{\fun1}$, which identifies $0$ with $\fun{nil}$ and $\fun{succ}(n)$ with $\fun{cons}(\fstar,n)$, this yields a path
    \[
        q\colon
        \dep{\N,(0,\fun{succ})}
        \eq{\fun L_{\fun1}\fun{Alg}}
        \dep{\fun L_{\fun1},(\fun{nil},\fun{cons})}.
    \]
    If we define the type family $Q\colon\fun L_{\fun1}\fun{Alg}\to\univ$ by $Q(\mathcal{A})\defeq\fun{isHinit}_{\fun L_{\fun1}}(\mathcal{A})$, then
    \[
        \fun{transport}^Q(q^{-1}) \colon 
        \fun{isHinit}_{\fun L_{\fun1}}(\dep{\fun L_{\fun1},(\fun{nil},\fun{cons})}) 
        \to 
        \fun{isHinit}_{\fun L_{\fun1}}(\dep{\N,(0,\fun{succ})}).
    \]
    Applying this function to the fact that $\dep{\fun L_{\fun1},(\fun{nil},\fun{cons})}$ is $h$-initial provides the desired inhabitant.
\end{proof}

\begin{defn}
    Let $A\colon\univ$ be a type of labels and $B\colon A\to\univ$ an arity family. A \textsl{polynomial functor} (in normal form) is the type family
    \[
        P(X)\defeq\sum_{a\colon A}B(a)\to X
    \]
    for any $X\colon\univ$.
\end{defn}

\begin{rem}
    The notion of polynomial functor can be seen as a generalization of \eqref{eq:P(W)->W}. The polynomial aspect of $P(X)$ becomes clear when $B(a)\to X$ is written as $X^{B(a)}$.
\end{rem}

\begin{defn} 
    The type of \textsl{$\fun W$-algebras} is
    \[
        \fun W\fun{Alg}(A,B)
            \defeq\sum_{C\colon\univ}\;
                \prod_{a\colon A}(B(a)\to C)\to C.
    \]
    Its terms are dependent pairs $\dep{C,s_C}$, where $s_C\colon\prod_{a\colon A}(B(a)\to C)\to C$ is the \textsl{structure map} of the $\fun W$-algebra.
\end{defn}

\begin{rem}
    A $\fun W$-algebra $\dep{C,s_C}$ admits an \textsl{uncurried variant} $\dep{C,\bar{s}_C}$, where
    \[
        \bar{s}_C\defeq\fun{rec}_{P(C)}(C,s_C)
            \colon P(C)\to C.
    \]
    We will refer to this form as a \textsl{$P$-algebra}.
    
    These two views are interchangeable since $s_C$ can be recovered from $\bar{s}_C$ by the computation rule of the $\Sigma$-type.
\end{rem}

\needspace{2\baselineskip}
\begin{defns}\ 
    \begin{enumerate}[-]
        \item Given $P$-algebras $\dep{C,\bar{s}_C}$ and $\dep{D,\bar{s}_D}$, any map $f\colon C\to D$ induces a function $Pf\colon P(C)\to P(D)$ defined by
        \[
            Pf(\dep{a,\phi})\defeq\dep{a,f\circ\phi},
        \]
        for $a\colon A$ and $\phi\colon B(a)\to C$.
        
        \item A \textsl{$P$-morphism} from a $P$-algebra $\dep{C,\bar{s}_C}$ to another $\dep{D,\bar{s}_D}$ is a dependent pair $\dep{f,\eta_f}$, where $f\colon C\to D$ and $\eta_f\colon f\circ\bar{s}_C\sim\bar{s}_D\circ Pf$. This homotopy can be visualized in the square:
        \[
            \begin{tikzcd}[font=\small]
                P(C)\arrow[r,"Pf"]
                        \arrow[d,"\bar{s}_C"']
                        \arrow[dr,
                            phantom,
                            "\scriptstyle\eta_f"]
                    &P(D)
                        \arrow[d,"\bar{s}_D"]\\
                C
                        \arrow[r,"f"']
                    &D
            \end{tikzcd}
        \]
        Thus, the type of $P$-morphisms is
        \[
            \fun P\Hom(\dep{C,\bar{s}_C},\dep{D,\bar{s}_D})
                \defeq\sum_{f\colon C\to D}
                    f\circ\bar{s}_C\sim\bar{s}_D\circ Pf.
        \]
        Equivalently, using the curried variants $s_C$ and $s_D$, the homotopy condition is given by a family of identities, making the type of morphisms definitionally equivalent to
        \begin{equation}\label{eq:structure-coherence}
            \sum_{f\colon C\to D}\;
                \prod_{a\colon A}\;
                \prod_{\phi\colon B(a)\to C}
                    f(s_C(a,\phi))\eq Ds_D(a,f\circ\phi).
        \end{equation}
        We will denote the type of $\fun W$-morphisms by $\fun W\Hom(\dep{C,s_C},\dep{D,s_D})$, or $\fun W\Hom(C,D)$ for short, if the context permits.

        \item A $\fun W$-algebra $\dep{I,s_I}$ is \textsl{$h$-initial} if the type
        \[
            \fun{isHinit}_{\fun W}(A,B,\dep{I,s_I})
            \defeq
            \prod_{\dep{C,s_C}\colon\fun W\fun{Alg}(A,B)}
            \!\!\!\fun{isContr}
            \bigl(
                \fun W\Hom(I,C)
            \bigr)
        \]
        is inhabited.
    \end{enumerate}
\end{defns}

\begin{defn}
    For a $\fun W$-type $\fun W\defeq\fun W_{x\colon A}B(x)$, the \textsl{recursor} is the function
    \[
        \fun{rec}_{\fun W}\colon
            \prod_{C\colon\univ}\;
            \Bigl(
                \prod_{a\colon A}(B(a)\to C)\to C
            \Bigr)
            \to\fun W\to C.
    \]
    For any type $C\colon\univ$ and structure map $s_C\colon\prod_{a\colon A}(B(a)\to C)\to C$, the recursion map $\fun{rec}_{\fun W}(C, s_C)\colon\fun W\to C$ is defined inductively on constructors by the rule
    \[
        \fun{rec}_{\fun W}(C,s_C)(\fun{sup}(a,\phi))
        \defeq
        s_C(a,\fun{rec}_{\fun W}(C,s_C)\circ\phi),
    \]
    for any $a\colon A$ and $\phi\colon B(a)\to\fun W$.
\end{defn}

\begin{ntn}\label{ntn:recursor-commutativity}
    For $\fun W\defeq\fun W_AB$, a carrier type $C\colon\univ$, and a curried structure map $s_C\colon\prod_{a\colon A}(B(a)\to C)\to C$, the term
    \[
        \fun{comp}_{\fun W}(C,s_C,a,\phi)
        \defeq\fun{refl}_{f(\fun{sup}(a,\phi))}
    \]
    denotes the witness for the computation rule \eqref{eq:rec_W-computation-rule}. Specifically, it is an inhabitant of the identity type
    \begin{equation}\label{eq:recursor-commutativity}
        \fun{rec}_{\fun W}(C,s_C)(\fun{sup}(a,\phi))
        \eq C
        s_C(a,\fun{rec}_{\fun W}(C,s_C)\circ\phi).
    \end{equation}
\end{ntn}

\begin{lem}\label{lem:transport-identity-family}
    Let $H, C\colon\univ$ be types and $L, R\colon H \to C$ functions. Given a path $p\colon h_1\eq{H} h_2$ and a path $q\colon L(h_1)\eq CR(h_1)$, we have
    \[
        \fun{transport}^{h\mapsto L(h)\eq CR(h)}(p,q) 
        \eq{L(h_2)\eq CR(h_2)}
        \fun{ap}_L(p)^{-1}\ct q\ct\fun{ap}_R(p).
    \]
\end{lem}

\begin{proof}
    First observe that, in fact, both sides of the target identity inhabit the type $L(h_2)\eq CR(h_2)$. Hence, we can apply path induction on $p$. It suffices to consider the case where $h_2\equiv h_1$ and $p\equiv\fun{refl}_{h_1}$. 
    
    In this case, the \lhs becomes definitionally $q$:
    \[
        \fun{transport}^{h \mapsto L(h) \eq{C} R(h)}(\fun{refl}_{h_1}, q) \equiv q.
    \]
    Since $\fun{ap}_L(\fun{refl}_{h_1})\equiv \fun{refl}_{L(h_1)}$ and $\fun{ap}_R(\fun{refl}_{h_1})\equiv \fun{refl}_{R(h_1)}$, the \rhs reduces to
    \[
        \fun{refl}_{L(h_1)}^{-1}
            \ct q
            \ct\fun{refl}_{R(h_1)}.
    \]
    The conclusion follows by the unit laws for path concatenation.
\end{proof}

\begin{lem}\label{lem:WHom-equality}
    Let $\dep{D,s_D}$ and $\dep{C,s_C}$ be two $\fun W$-algebras. Let $\dep{f,\eta_f}$ and $\dep{g,\eta_g}$ be two $\fun W$-morphisms from $\dep{D,s_D}$ to $\dep{C,s_C}$. The identity type
    \[
        \dep{f,\eta_f}\eq{\fun W\Hom(D,C)}\dep{g,\eta_g}
    \]
    is equivalent to the type of dependent pairs $\dep{e,s_e}$, where $e\colon f\sim g$ is a homotopy and $s_e$ is a family of paths assigning to each $a\colon A$ and $\phi\colon B(a)\to D$ a path
    \[
        s_e(a,\phi)\colon
            e(s_D(a,\phi))
            \eq{}
            \eta_f(a,\phi)
            \ct\fun{ap}_{s_C(a,\dummy)}(
                \fun{funext}(e\rwct\phi))
            \ct\eta_g(a,\phi)^{-1}
    \]
    in $f(s_D(a,\phi)) \eq{C} g(s_D(a,\phi))$. In other words, the path diagram
    \[
        \begin{tikzcd}[row sep=large, column sep=3.2cm]
            f(s_D(a,\phi))
                    \arrow[r,"{e(s_D(a,\phi))}"]
                    \arrow[d,"{\eta_f(a,\phi)}"']
                    \arrow[rd,
                        "{\scriptstyle s_e(a,\phi)}",
                        phantom]
                &g(s_D(a,\phi))
                    \arrow[d,gray,"{\eta_g(a,\phi)}"]\\
            s_C(a,f\circ\phi)
                    \arrow[r,
                        "{\fun{ap}_{s_C(a,\dummy)}
                            (\fun{funext}(e\rwct\phi))
                        }",
                        swap]
                &s_C(a,g\circ\phi)
                    \arrow[u,
                        bend left=50,
                        "{\eta_g(a,\phi)^{-1}}"]
        \end{tikzcd}
    \]
    commutes propositionally if, and only if, $\dep{f,\eta_f}\eq{\fun W\Hom(D,C)}\dep{g,\eta_g}$.
\end{lem}

\begin{proof}
    Equating the dependent pairs $\dep{f,\eta_f}$ and $\dep{g,\eta_g}$ in the $\Sigma$-type of $\fun W$-morphisms $\fun W\Hom(D,C)$ is equivalent to providing a path $p\colon f\eq{D\to C}g$ and a path $q$ witnessing the identity between the transport of $\eta_f$ along $p$ and $\eta_g$.
    
    By function extensionality, the path $p$ is equivalent to a homotopy $e\colon f\sim g$, giving $p\defeq\fun{funext}(e)$.
    
    For any $a\colon A$ and $\phi\colon B(a)\to D$, let us define the local terms
    \[
        L_{a,\phi}(h)\defeq h(s_D(a,\phi))
            \quad\text{and}\quad
        R_{a,\phi}(h)\defeq s_C(a,h\circ\phi).
    \]
    This allows us to define the dependent type family $E$ as
    \[
        E(h)\defeq
            \prod_{a\colon A}\;
            \prod_{\phi\colon B(a)\to D}
                L_{a,\phi}(h)\eq CR_{a,\phi}(h).
    \]
    According to \eqref{eq:structure-coherence}, we have $\eta_f\colon E(f)$ and $\eta_g\colon E(g)$, and the transport condition mentioned above translates to
    \[
        \fun{transport}^E(p,\eta_f)\eq{}\eta_g.
    \]
    By applying Proposition~\ref{prop:transport-pointwise-evaluation} to the outer product over $A$, we obtain
    \[
        \fun{transport}^E(p,\eta_f)
        \eq{}
        \lambda a.\,\fun{transport}^{
            \lambda h.\,\prod_{\phi\colon B(a)\to D}
                L_{a,\phi}(h)\eq CR_{a,\phi}(h)
            }
            (p,\eta_f(a)).
    \]
    Applying the proposition a second time to the inner product over $B(a)\to D$ yields
    \[
        \fun{transport}^E(p,\eta_f)\eq{}
        \lambda a.\,\lambda \phi.\, 
            \fun{transport}^{
                \lambda h.\,L_{a,\phi}(h)
                \eq CR_{a,\phi}(h)
            }
            (p,\eta_f(a,\phi)).
    \]
    By Lemma~\ref{lem:transport-identity-family}, we obtain
    \[\tag{$\ast$}
        \fun{transport}^{h \mapsto L_{a,\phi}(h) \eq{C} R_{a,\phi}(h)}(p, \eta_f(a,\phi)) 
        \eq{} \fun{ap}_{L_{a,\phi}}(p)^{-1} \ct \eta_f(a,\phi) \ct \fun{ap}_{R_{a,\phi}}(p).
    \]
    \needspace{2\baselineskip}
    Computing the action of $\fun{ap}_{L_{a,\phi}}$ and $\fun{ap}_{R_{a,\phi}}$ for $p\defeq\fun{funext}(e)$ gives:
    \begin{enumerate}[a),font=\upshape]
        \item $\fun{ap}_{L_{a,\phi}}(\fun{funext}(e))\eq{} e(s_D(a,\phi))$.
        \item $\fun{ap}_{R_{a,\phi}}(\fun{funext}(e))\eq{} \fun{ap}_{s_C(a,\dummy)}(\fun{funext}(e\rwct\phi))$.
    \end{enumerate}
    Substituting these into the transport equation ($\ast$) yields
    \[
        e(s_D(a,\phi))^{-1}
        \ct\eta_f(a,\phi)
        \ct\fun{ap}_{s_C(a,\dummy)}
            (
                \fun{funext}(e\rwct\phi)
            ).
    \]
    The coherence condition requires this transported path to equal $\eta_g(a,\phi)$:
    \[
        e(s_D(a,\phi))^{-1} \ct \eta_f(a,\phi) \ct \fun{ap}_{s_C(a,\dummy)}(\fun{funext}(e\rwct\phi)) \eq{} \eta_g(a,\phi).
    \]
    Concatenating on the left by $e(s_D(a,\phi))$ and on the right by $\eta_g(a,\phi)^{-1}$ isolates $e(s_D(a,\phi))$, yielding the required coherence condition:
    \[
        e(s_D(a,\phi))
        \eq{}
        \eta_f(a,\phi)
            \ct\fun{ap}_{s_C(a,\dummy)}
                (
                    \fun{funext}(e\rwct\phi)
                )
            \ct\eta_g(a,\phi)^{-1}.
    \]
    This completes the proof of the \textit{if} part.
    
    For the \textit{only if} part, note that the equality $\dep{f,\eta_f}\eq{}\dep{g,\eta_g}$ means the existence of paths
    \begin{align*}
        p&\colon f\eq{D\to C}g,\\
        q&\colon
            \fun{transport}^E(p,\eta_f)
            \eq{E(g)}
            \eta_g.
    \end{align*}
    Thus, we can apply $\fun{happly}$ to $q$ to translate this global equality into the local path:
    \[
        \fun{happly}(q,a)\colon
            \fun{transport}^E(p,\eta_f)(a)
            \eq{B(a)\to D}
            \eta_g(a).
    \]
    By applying $\fun{happly}$ a second time we obtain
    \[
        \fun{happly}(\fun{happly}(q,a),\phi)\colon
        \fun{transport}^E(p,\eta_f)(a,\phi)
        \eq{}
        \eta_g(a,\phi).
    \]
    According to Proposition~\ref{prop:transport-pointwise-evaluation} and Lemma~\ref{lem:transport-identity-family}, the left-hand side reduces point-wise. Thus, this double application of $\fun{happly}$ yields the local path 
    \[
        e(s_D(a,\phi))^{-1}
        \ct\eta_f(a,\phi)
        \ct\fun{ap}_{s_C(a,\dummy)}
            (
                \fun{funext}(e\rwct\phi)
            )\eq{}\eta_g(a,\phi),
    \]
    which we can use to derive $s_e(a,\phi)$.
\end{proof}

\begin{thm}
    For any type $A\colon\univ$ and type family $B\colon A\to\univ$, the $\fun W$-algebra $\dep{\fun W_AB,\fun{sup}}$ is $h$-initial.
\end{thm}

\begin{proof}
    Let $\fun W\defeq\fun W_AB$. We must show that for any $\fun W$-algebra $\dep{C,s_C}$, the type $\fun W\Hom(\fun W,C)$ is contractible. 
    
    For the center of contraction $\dep{f,\eta_f}$, let $f\colon\fun W\to C$ be defined by \eqref{eq:rec_W}, i.e.,
    \[
        f\defeq\fun{rec}_{\fun W}(C,s_C).
    \]
    The homotopy condition \eqref{eq:structure-coherence} for $\eta_f$ requires that, for $a\colon A$ and $\phi\colon B(a)\to\fun W$, we have
    \[
        f(\fun{sup}(a,\phi))\eq Cs_C(a,f\circ\phi).
    \]
    This equality is exactly the computation rule \eqref{eq:recursor-commutativity} for the $\fun W$-recursor. Thus, using Notation~\ref{ntn:recursor-commutativity}, we define $\eta_f(a,\phi)\defeq\fun{comp}_{\fun W}(C,s_C,a,\phi)$.

    To prove contractibility, it remains to show that for any other $\fun W$-morphism $\dep{g,\eta_g}$, there is an identity $\dep{f,\eta_f}\eq{\fun W\Hom(\fun W,C)}\dep{g,\eta_g}$. 
    
    By Lemma~\ref{lem:WHom-equality}, providing such an identity is equivalent to providing a dependent pair $\dep{e,s_e}$, where $e\colon f\sim g$ and $s_e(a,\phi)$ witnesses the equality
    \[
        e(\fun{sup}(a,\phi))\eq{}
        \eta_f(a,\phi)
        \ct\fun{ap}_{s_C(a,\dummy)}
            (\fun{funext}(e\rwct\phi))
        \ct \eta_g(a,\phi)^{-1}.
    \]
    To construct the homotopy $e$, we proceed by $\fun W$-induction for the type family $E\colon\fun W\to\univ$ defined by
    \[
        E(w)\defeq f(w)\eq Cg(w).
    \]
    According to \eqref{eq:ind_W-step-function}, given
    \[
        \epsilon\colon\prod_{b\colon B(a)}
            f(\phi(b))\eq Cg(\phi(b)),
    \]
    we need to provide a path in $f(\fun{sup}(a,\phi))\eq Cg(\fun{sup}(a,\phi))$. We define
    \[
        e_s(a,\phi,\epsilon)\defeq
            \eta_f(a,\phi)
            \ct\fun{ap}_{s_C(a,\dummy)}
                (\fun{funext}(\epsilon))
            \ct\eta_g(a,\phi)^{-1}.
    \]
    By the computation rule \eqref{eq:ind_W-computation-rule}, the resulting homotopy $e\colon f\sim g$ satisfies
    \[
        e(\fun{sup}(a,\phi))\eq{}e_s(a,\phi,e\rwct\phi),
    \]
    since $e\rwct\phi\equiv e\circ\phi$.
    
    Thus, the required coherence condition $s_e$ of Lemma~\ref{lem:WHom-equality} is satisfied, yielding the required identity $\dep{f,\eta_f}\eq{\fun W\Hom(\fun W,C)}\dep{g,\eta_g}$.
    %
    
\end{proof}
